\documentclass[a4paper,11pt]{report}
% The legacy CTAN style is bundled for projects that use USCthesis. It is
% loaded here to expose its report metadata helpers while the portable example
% uses the modern report class for standalone compilation.
\usepackage{techreport}
\usepackage[T1]{fontenc}
\usepackage{lmodern}
\usepackage{booktabs}
\usepackage{hyperref}
\hypersetup{colorlinks=true,linkcolor=blue,urlcolor=blue}
\reportnumber{TR-2026-001}
\address{Research and Operations Group}
\acknowledgment{Prepared as a neutral example for technical-report authors.}

\begin{document}
\begin{titlepage}
  \centering
  \vspace*{2cm}
  {\Large TECHNICAL REPORT\par}
  \vspace{1.5cm}
  {\huge\bfseries Portable Technical Report Starter\par}
  \vspace{1cm}
  {\large TR-2026-001\par}
  \vfill
  Research and Operations Group\par
  August 2026
\end{titlepage}

\chapter{Purpose}
This document demonstrates a portable report layout alongside the metadata
commands provided by the legacy CTAN \texttt{techreport} style. The bundled
style was originally designed as an option for the USC thesis class; projects
using that class can adapt this starter by changing the document class and
invoking its title-page hooks.

\chapter{Status and Findings}
The current phase is complete enough for review. The following compact table
shows the information typically carried by a technical report.

\begin{table}[ht]
  \centering
  \caption{Illustrative status summary}
  \begin{tabular}{lll}
    \toprule
    Area & Status & Next action \\
    \midrule
    Scope & Confirmed & Publish baseline \\
    Evidence & Reviewed & Archive sources \\
    Delivery & On track & Schedule handoff \\
    \bottomrule
  \end{tabular}
\end{table}

\chapter{Recommendations}
Record the report number, acknowledgement, address, and title-page policy in
the project configuration before using this document for a real deliverable.

\end{document}
